\subsection*{119. Study the Distribution of Learners, Not Only the Mean}

A cohort mean can conceal essential realities. Some learners enter with strong foundations and gain little. Others struggle but improve greatly. Some leave early. Some fail to demonstrate an essential competence while scoring highly on easier tasks.

Track starting-state distributions, important error classes, participation, mastery, assistance levels, and the number of people observed at each stage. An average exit score of 80\% is difficult to interpret if half the cohort did not take the assessment or if the test was rehearsed.

Avoid claiming precise value added from a small pre/post sample. Chapter 13 will examine stronger designs. For now, learn to keep the \textbf{denominator, assistance policy, and alternative explanations visible}.

\subsection*{120. Test the Pedagogy Before Industrializing the Content}

A founder can produce a beautiful library of 100 videos only to discover that learners repeatedly misunderstand the core idea or lack time for the promised work. The cheaper route is to teach a small representative sample, inspect authentic attempts, and revise the sequence before scaling production.

A good early experiment might include 1 strong reading, 1 concise explanation, 2 worked cases, 3 independent variants, feedback, a retry, and a later transfer task. Its success criterion is a defensible learning gain on the selected skill, not the attractiveness of the studio recording.

Production quality is important when it reduces confusion, improves accessibility, or supports reuse. It should be purchased in proportion to proven educational need.

\subsection*{121. Measure the Scarce Resources That Actually Create the Change}

For a technical cohort, time budgets should separate explanations, learner practice, human review, administration, admissions, curriculum updates, support, and remedial work. Track actual rather than planned effort. If a mentor spends 2 hours rescuing every student's project, the cost and learner dependence matter.

For operational reasoning, define:
\begin{equation}
\text{Review hours} = \sum_{i=1}^{N}\sum_{j=1}^{K_i} \frac{M_{i,j}}{60},
\end{equation}
where $M_{i,j}$ represents the minutes spent reviewing attempt $j$ of learner $i$, summed across all $N$ learners and their respective $K_i$ attempts.

The expression is simple but discourages vague promises. Compare review intensity against the \textbf{quality of subsequent independent work}, not just completion or satisfaction. Distinguish cash expenses from normalized labor costs and allow contingency for uneven learner needs.

\subsection*{122. Separate 3 Tests: Did They Learn, Did They Value It, Can We Repeat It?}

These are the 3 questions of educational product discipline:

\textbf{Learning:} Did representative participants demonstrate a consequential improvement, with honest limitations? \textbf{Value:} Would suitable buyers accept the money, time, and effort required for this experience compared with alternatives? \textbf{Repeatability:} Can the institution recruit, support, and teach subsequent comparable groups without unsustainable founder labor or deteriorating standards?

A program can pass 1 and fail the others. This is not a contradiction. It tells the founder which system needs repair. The logical next action may be to change the buyer, reduce scope, redesign feedback, alter the duration, or discontinue the offer. It is rarely to add another inspirational slogan.

\subsection*{123. Build a Teacher-Learning Loop}

After each delivery, record the intended capability, common misconceptions, observed attempts, which interventions were used, what changed after feedback, and which tasks remain unreliable. Invite faculty to propose alternative explanations for the failures.

For example: \emph{Learners could recite the need for temporal validation but still mishandled repeated customers.} Perhaps the lesson overemphasized time while underexplaining entity dependence. Add contrasting cases, then test the revised explanation on an unseen problem in the next delivery.

A serious institution should be \textbf{scientifically curious about its own pedagogy}. It may not have the sample size for decisive randomized experiments immediately, but it can maintain honest, cumulative evidence and avoid pretending that satisfaction ratings prove mastery.

\subsection*{124. Know When to Stop, When to Narrow, and When to Deepen}

A course may fail because its target capability is too broad, the learners lack prerequisites, delivery requires unaffordable expert labor, or the intended customer does not value the outcome enough to pay. More content does not automatically solve any of these.

\textbf{Narrow} when 1 sub-capability produces clear value while others dilute effort. \textbf{Deepen} when learners can complete prepared tasks but fail to reason in unfamiliar situations. \textbf{Redesign} when the feedback loop is weak. \textbf{Stop or change market} when a sound educational product cannot be sold or sustainably delivered to the intended buyer under acceptable conditions.

These decisions preserve the founder's long-term mission by refusing to finance it with a failing short-term mechanism.

\subsection*{125. The Learning Founder's Final Test}

Before approving a new program, the founder should be able to answer:

\begin{quote}
What will learners be able to do? Why does it matter? What do they already know? Which cognitive and behavioral changes are necessary? What attempts will create those changes? Who or what will detect important errors? How will learners act on feedback? Where will the knowledge be retrieved later? How will it be used under different conditions? How will we know what happened, and can we afford to repeat the experience?
\end{quote}

A founder who can answer these questions has not guaranteed learning. They have done something more valuable than publishing a syllabus: constructed a falsifiable theory of transformation and the beginnings of an operating system that can improve itself.
